\section{The lagging paradox, formally}\label{sec:lagging}

A recurring observation in applied work on deep price forecasting is that networks trained with
squared error on price \emph{levels} produce forecasts that ``lag'' the realised path: the
predicted curve looks like the actual curve shifted one day to the right, and directional accuracy
is close to or below one half even when the price-level RMSE looks impressive. We call this the
\emph{lagging paradox}. This section traces the lag to the loss function, whatever the architecture.
It derives a sharp upper bound on what any price-level forecaster can gain over the random walk and
shows that a widely used directional metric measures the sign persistence of returns, which says
nothing about forecasting skill. The results motivate the distributional formulation in
\cref{sec:model}.

\subsection{Setting}
Let $P_t>0$ be the price at the close of day $t$, $r_{t+1}=\log P_{t+1}-\log P_t$ the next log
return and $\F_t$ the information at the close of day $t$. Write
$M_t(s)=\E[e^{s r_{t+1}}\mid\F_t]$ for the conditional moment generating function, assumed finite
for $s\in\{1,2\}$, and $\mu_t=\E[r_{t+1}\mid\F_t]$, $\sigma_t^2=\Var(r_{t+1}\mid\F_t)$.

\subsection{The MSE-optimal price forecast is a near-random walk}

\begin{proposition}[Optimal level forecast]\label{prop:optimal-level}
Among $\F_t$-measurable forecasts $f_t$ of $P_{t+1}$, the conditional mean squared error
$\E[(P_{t+1}-f_t)^2\mid\F_t]$ is minimised by $f_t^\ast=P_t\,M_t(1)$, with minimum value
$P_t^2\{M_t(2)-M_t(1)^2\}$. The random-walk forecast $f_t^{\mathrm{RW}}=P_t$ has conditional MSE
$P_t^2\{M_t(2)-2M_t(1)+1\}$, and
\begin{equation}\label{eq:mse-gap}
  \E[(P_{t+1}-P_t)^2\mid\F_t]-\E[(P_{t+1}-f_t^\ast)^2\mid\F_t]=P_t^2\{M_t(1)-1\}^2 .
\end{equation}
\end{proposition}

\noindent Under conditional normality $M_t(1)=\exp(\mu_t+\sigma_t^2/2)$, so the optimal forecast
moves away from the random walk by the factor $\exp(\mu_t+\sigma_t^2/2)-1\approx m_t$,
$m_t\coloneqq\mu_t+\sigma_t^2/2$. For daily equity returns $m_t$ is of order $10^{-4}$, so
$\Cov(f^\ast_t,P_t)/\{\Var(f^\ast_t)\Var(P_t)\}^{1/2}$ is indistinguishable from one: \emph{the
optimal forecast lags by construction.}

\begin{proposition}[Upper bound on price-level gains]\label{prop:bound}
Let $s_t=m_t/\sigma_t$ be the conditional drift-to-noise ratio. Under conditional normality and
to second order in $(\mu_t,\sigma_t)$, the relative reduction in conditional MSE that any
$\F_t$-measurable forecast can achieve over the random walk satisfies
\begin{equation}\label{eq:bound}
  \frac{\E[(P_{t+1}-P_t)^2\mid\F_t]-\min_f\E[(P_{t+1}-f)^2\mid\F_t]}{\E[(P_{t+1}-P_t)^2\mid\F_t]}
  \;=\;\frac{s_t^2}{1+s_t^2}+O(\sigma_t^3).
\end{equation}
Consequently the unconditional relative RMSE gain is at most
$1-\{1-\E[w_t s_t^2/(1+s_t^2)]\}^{1/2}\approx \tfrac12\E[w_t s_t^2]$ with weights
$w_t=P_t^2\sigma_t^2/\E[P_t^2\sigma_t^2]$.
\end{proposition}

\noindent A daily drift-to-noise ratio of $|s_t|=0.1$ (an extreme degree of predictability,
corresponding to an annualised conditional Sharpe ratio of about $1.6$ on \emph{every} day) bounds
the achievable price-RMSE improvement at roughly $0.5\%$. Reported RMSE improvements of tens of
percent on daily index levels are therefore incompatible with any information set available at
$t$ and indicate look-ahead in the evaluation. In \cref{sec:results-china} we confirm empirically that a
60-lag linear model fitted by least squares to price levels converges to the random walk on all
nineteen assets in our sample (sum of coefficients $\approx 1$) and loses between $0.6\%$ and
$7.8\%$ in RMSE out of sample, as the bound and estimation noise predict
(\cref{tab:lagging}).

\subsection{Differenced directional accuracy measures sign persistence}

Directional accuracy is often computed from consecutive forecasts as
$\mathrm{DA}^{\Delta}=\frac1n\sum_t\1\{\operatorname{sign}(f_{t+1}-f_{t})=\operatorname{sign}(P_{t+1}-P_t)\}$,
i.e.\ comparing the change in the forecast path with the change in the price path.

\begin{proposition}[What $\mathrm{DA}^{\Delta}$ measures]\label{prop:da}
For the random-walk forecast $f_{t+1}=P_t$,
$\mathrm{DA}^{\Delta}=\frac1n\sum_t\1\{\operatorname{sign} r_t=\operatorname{sign} r_{t+1}\}$, the empirical
sign persistence of returns, whose expectation is
$\tfrac12+\tfrac12\E[\operatorname{sign}(r_t)\operatorname{sign}(r_{t+1})]$ when $P(r_t=0)=0$. For any forecast
with $f_{t+1}=P_t+\varepsilon_{t+1}$ and $|\varepsilon_{t+1}-\varepsilon_t|<|r_t|$ the same identity
holds term by term.
\end{proposition}

\noindent Whenever the forecaster lags, $\mathrm{DA}^{\Delta}$ is therefore a statistic of the data
alone. The skill-relevant quantity compares the forecast with the \emph{current} price,
$\mathrm{DA}=\frac1n\sum_t\1\{\operatorname{sign}(f_{t+1}-P_t)=\operatorname{sign}(r_{t+1})\}$, which is undefined for the
random walk and is what we report throughout.

\subsection{Why the conditional distribution is the right target}

Under 0--1 loss the optimal directional call is $\1\{P(r_{t+1}>0\mid\F_t)>1/2\}$, i.e.\ the sign of
the conditional \emph{median}, not of the mean. For a Gaussian the two coincide; for a mixture they
need not.

\begin{example}\label{ex:median}
Let $r_{t+1}\mid\F_t\sim 0.8\,\N(0.3,1)+0.2\,\N(-1.0,3^2)$. The mean is $0.04>0$, but
$P(r_{t+1}>0)=0.8\,\Phi(0.3)+0.2\,\Phi(-1/3)=0.568$; changing the second component to
$\N(-1.5,3^2)$ gives mean $-0.06<0$ while $P(r_{t+1}>0)=0.556>1/2$. A point forecaster trained by
squared error calls ``down''; the distributional forecaster calls ``up'' and also reports that a
fifth of the probability mass sits in a high-variance regime.
\end{example}

\noindent These results organise the rest of the paper. First, price-level losses
cannot reveal predictive skill; the target must be the return. Second, a forecaster should be judged
by proper scoring rules for the whole predictive distribution, since direction, volatility and tail
risk are functionals of it. Third, if regimes matter, the predictive distribution is naturally a
mixture, and a mixture whose components and weights are \emph{readable} is the white-box object
we construct next.
